\section{Principal mathematical conclusions}\label{sec:conclusions}
Each conclusion cites the result on which it rests; broad readings are labelled as interpretation.
\subsection{Operational characterisation of the UNNS Substrate}\label{sec:characterisation}
The preceding results support an operational characterisation of the UNNS Substrate as a rule-governed generative arithmetic architecture. Its mathematical organisation is not exhausted by numerical states, individual generating histories, occurrence inventories, primitive arithmetic relations, or realised geometric complexes considered separately. It is expressed through the admissible transformations and structural relationships among these objects.

The generative framework may be represented schematically by
\[
\mathcal U=(S,\mathcal O,\mathrm{Dom},\mathcal R,\mathsf{Adm}),
\]
where $S$ is a set of states, $\mathcal O$ a set of operators, $\mathrm{Dom}$ the assignment of a domain to each operator, $\mathcal R$ the set of admissible routes, and $\mathsf{Adm}$ the conditions governing admissibility. In the arithmetic constructions studied here, $S$ consists of the positive integers and rationals, $\mathcal O$ of the maps $T_N$ and $C$, $\mathrm{Dom}$ of the conditions $N\mid m$ and $D(x)\neq\emptyset$, $\mathcal R$ of histories and collective orbits, and $\mathsf{Adm}$ of the divisibility conditions on squares. The tuple is a notation for organising what has been proved; it is not an axiom system from which the results are derived. The occurrence and geometric structures studied here ($D$, the lattice of relations, the primitive system $\mathrm{Gr}_q$, the complexes $K(D)$ and their Betti numbers) are associated constructions, not additional primitive components assumed without definition.

The results establish the following characteristic properties of this architecture; each is a restatement of the cited result for the stated domain.
\begin{center}\footnotesize\setlength{\tabcolsep}{4pt}\begin{tabular}{@{}p{7.2cm}p{4.6cm}@{}}\toprule
Property & Support\\\midrule
Numerical endpoints need not identify their generative histories; the endpoint retains only the valuation displacement. & \cref{prop:P1}, \cref{prop:P5}\\
Occurrence inventories determine the realised move complexes, but not the complete primitive relation systems. & \cref{prop:P2}(b),(c); \Src{III}{thm:R3T1a}, \Src{III}{thm:R3T1}\\
Arithmetic relations can remain dormant until further occurrences permit their geometric realisation. & \cref{prop:P3}; \Src{I}{thm:P2T2}\\
Global connectivity of the history space is governed by arithmetic of the source, not by local swaps alone. & \cref{prop:P4}\\
Under the collective map, occurrence sets grow, connectivity only coarsens, and each edge exists during an arithmetically determined interval. & \cref{prop:P3}(b), \cref{prop:P6}(a)--(c)\\
Event steps follow exact arithmetic laws, and at graph level the change of the first two Betti numbers follows from local counts and prime supports; local counts alone do not determine it, and beyond graph level the determination is open. & \cref{prop:P7}; \Src{III}{thm:sieve}\\
Scaling parameters do not determine topology (finite-exact). & \cref{prop:P8}\\
The two rules are related at integers and occurrence labels, not by a one-step identity correspondence. & \cref{prop:P9}; \cref{fig:global}\\
\bottomrule\end{tabular}\end{center}

These properties show that generation, arithmetic possibility, geometric realisation and structural evolution are distinct but mathematically related aspects of the framework. Their relationships are governed by the proved factorizations, admissibility conditions, transition laws and obstructions of \cref{sec:numerical,sec:occurrence,sec:realization,sec:history,sec:endpoint,sec:evolution,sec:prediction,sec:recurrence,sec:correspondence}, collected in \cref{fig:global}.

Consequently, the UNNS Substrate is characterised operationally not by a single geometric object, but by a system of admissible generative mechanisms and the layered structures they produce and transform.

\emph{Scope.} This characterisation is supported by the results established for the specified arithmetic constructions. It does not assert that every component derives from a unique minimal axiom system, that all generative operators are dynamically equivalent, or that the mathematical framework has been identified with a physical substrate. Establishing a fully intrinsic characterisation remains a separate mathematical problem (\cref{sec:intrinsic}).

\subsection{Mathematical implications}\label{sec:implications}
Each conclusion cites the result on which it rests; broad readings are labelled as interpretation.
\begin{enumerate}[label=\textbf{C\arabic*.},itemsep=3pt]
\item \textbf{Generation.} Arithmetic conditions decide which steps and which histories exist: a step $m\xrightarrow{N}t$ exists iff $d=N+1$ satisfies $d^2\mid t$ (\cref{prop:P2}(a)), and legal swaps are decided by divisibility at each state (\cref{prop:P4}).
\item \textbf{Ancestry.} The numerical endpoint determines only the valuation displacement of the history (\cref{prop:P1}); equal occurrence sets are compatible with different numbers of histories and of swap classes (\cref{prop:P5}).
\item \textbf{Geometry.} The realised geometry of a state is a function of its occurrence set (\Src{III}{thm:R3T1a}); the primitive relation system contains information that the occurrence set does not (\cref{prop:P2}(c)).
\item \textbf{Transport.} Preservation under lattice inclusion holds in rank two (\Src{I}{thm:planar}) and fails in rank three under the hypotheses of \Src{I}{thm:C1}; abstract and realised loss differ (\Src{I}{prop:topology}).
\item \textbf{Recombination.} Local legal swaps do not force a connected history space: the class count of $\{2,3,14\}$ is $0$, $1$ or $2$ according to arithmetic of the source and is not monotone in divisibility (\cref{prop:P4}).
\item \textbf{Evolution.} Collective steps enlarge the occurrence set, only coarsen the connectivity relation, give each edge an interval of existence, and leave higher Betti numbers unconstrained by monotonicity (\cref{prop:P3,prop:P6}).
\item \textbf{Prediction.} Arithmetic congruences give the event step; supports and local counts give $\Delta\beta_0$ and, at graph level, $\Delta\beta_1$; positive-degree homology needs attachment data (\cref{prop:P7}, \Src{III}{thm:sieve}, \Src{III}{thm:zigzag}). Scaling data do not determine topology (\cref{prop:P8}).
\item \textbf{Compatibility.} The one-step identity correspondence of the two rules is obstructed (a non-identity intertwining is open); the overlap on integers admits a translation formula; the other levels are obstructed or open (\cref{prop:P9}).
\end{enumerate}
\paragraph{Answer to the governing question.} \emph{Interpretation, supported by C1--C8 and not a theorem:} in the part of the theory that is established, geometry is formed from arithmetic by two successive forgetful maps---from histories to endpoints, from states to occurrence sets---and the formation of geometry is stable in a precise sense (occurrence sets only grow; connectivity only coarsens) but not rigid (edges and higher homology are created and destroyed). Where the two generating rules meet, they meet in integers and in occurrence labels, and the exact obstructions show that nothing richer follows from the labels alone.

\subsection{Limitations}\label{sec:limitations}
The conclusions concern the specified arithmetic constructions and the finite-exact ranges stated in the propositions. Verification inherited from Papers~I--III is not re-run here (NOT\_RERUN). The collective rule is studied on its recorded domains and the generative rule on integer sources; the extension of the generative rule to rational sources is open (O8). No claim is made that the operators $T_N$ and $C$ belong to one dynamical system, that the associated constructions are independent of one another beyond the stated obstructions, or that the characterisation applies outside the constructions named. Constant C3 is unclaimed.
